\sscp{Secondary Papuan incursion(s): east-to-west}{15-06-04}
As described in \srf{02-01},
there is considerable archaeological evidence that the Timor
region had human habitation by at least 42,000 BP \citep{OConnor2015}.
We do not know what languages these early inhabitants spoke.
While some people appear to have travelled onward to the New
Guinea region and to Australia, the archaeological record shows that others
stayed in the Timor region for thousands of years.

By the time the Austronesian-speaking peoples arrived around
4,000--3,500 BP, the non-Austronesian languages spoken in the region seem to have
been predominantly typologically SOV,
with preposed possessors, post-verbal negation,
and an alienable-inalienable distinction in their possessive systems.
These features are still widely found in the Papuan \tsc{Timor-Alor-Pantar}
languages in the Timor region,
as well as in some \tsc{North Halmahera} languages,
and other Papuan languages of
the Bird’s Head of Papua.

As discussed in \srf{02-01},
the idea of \tsc{Timor-Alor-Pantar} (TAP) languages linking to
Papua through Proto-Greater West Bomberai has been gaining increasing
support \citep{UsherSchapper2022}.
The accumulating typological and lexical evidence for such an
east-to-west link also parallels the westward movement away from
the New Guinea region of important food crops
and marsupials that we have discussed in \srf{02-03-01-07} and \crf{ch:Mar}.
These occurred several thousand years before the arrival of the
Austronesian speaking peoples \citep{DenhamDonohue2009,DonohueDenham2009,
DonohueDenham2010,FlanneryNitihaminoto1995,Schapper2011}.
The east-to-west spread of the term {\#}muku ‘banana’ from the
New Guinea region along the southern islands all the way to western
Flores, but not in most of central Maluku,
is only one example.